\documentclass[english,11pt,a4paper]{article}
\usepackage[margin=1in]{geometry}
\usepackage{fontspec}
\setmainfont{lmroman10-regular.otf}[BoldFont=lmroman10-bold.otf,ItalicFont=lmroman10-italic.otf,BoldItalicFont=lmroman10-bolditalic.otf]
\setsansfont{lmsans10-regular.otf}[BoldFont=lmsans10-bold.otf]
\setmonofont{lmmono10-regular.otf}[Scale=MatchLowercase]
\newfontfamily\LeanSymbols{seguisym.ttf}[Path=typeset_runtime/]
\newlength{\LeanCharWidth}
\usepackage{amsmath,amssymb}
\usepackage[english]{babel}
\usepackage{microtype,parskip}
\usepackage{graphicx,float,caption}
\usepackage{longtable,booktabs,array,calc}
\usepackage{xcolor}
\usepackage{xurl}
\usepackage{hyperref,bookmark}
\usepackage{etoolbox,needspace}
\setcounter{secnumdepth}{-\maxdimen}
\setlength{\emergencystretch}{3em}
\providecommand{\tightlist}{\setlength{\itemsep}{0pt}\setlength{\parskip}{0pt}}
\newcounter{none}
\AtBeginEnvironment{longtable}{\Needspace{8\baselineskip}\small}
\pretocmd{\section}{\Needspace{5\baselineskip}}{}{}
\pretocmd{\subsection}{\Needspace{4\baselineskip}}{}{}
\pretocmd{\subsubsection}{\Needspace{4\baselineskip}}{}{}
\captionsetup{font=small,labelfont=bf}
\urlstyle{same}
\providecommand{\pandocbounded}[1]{#1}
\hypersetup{pdftitle={Paper III: formalization and reproducibility},pdfauthor={Juan Pablo Traverso Gianini},pdflang={en-US},colorlinks=true,linkcolor=blue,urlcolor=blue,citecolor=blue}
\title{Paper III: formalization and reproducibility}
\author{Juan Pablo Traverso Gianini\\\small Independent researcher, Santiago, Chile\\\small \href{mailto:jtraverso@gmail.com}{jtraverso@gmail.com}\\\small ORCID: 0009-0003-6068-4096}
\date{}
\begin{document}
\raggedbottom
\maketitle
\textbf{Paper III · Formal companion v1.6}\\
19 October 2026

This document accompanies the mathematical article. It preserves the
formal declarations, scope qualifications and computational evidence
previously included in the article. The source freeze and its recorded
verification are unchanged. No Lean/Lake build was run for this edition.
Historical section and table numbers in the records below refer to the
preceding article; they are retained as locators, not renumbered as new
verification results. Original references {[}1{]}, {[}2{]}, etc. are
listed at the end.

\href{https://github.com/jtraverso/erdos-81-chordal-clique-partitions/tree/eb335845bfa23375d5040ea31fc4f4b366346287/preprints/PAPER_III/05_formalization}{Frozen
sources and historical evidence / Fuentes congeladas y evidencias
históricas} ·
\href{https://github.com/jtraverso/erdos-81-chordal-clique-partitions/tree/eb335845bfa23375d5040ea31fc4f4b366346287/preprints/PAPER_III/01_manuscript}{Previous
article / Artículo anterior}

\subsection{11.6 Formal-verification perimeter and current
status}\label{formal-verification-perimeter-and-current-status}

The formalization discussion is confined to this section and the release
record in Section 13. The current formal snapshot separates the
unconditional final assembly from its downstream interfaces. The public
aggregate root \nolinkurl{PaperIII} imports both
\nolinkurl{PaperIII.Theorem_1_1_Final} and \nolinkurl{PaperIII.PublicAPI}.
The final module discharges the two interfaces used by the paper and
exports Theorem 1.1 and Corollary 1.2.
\nolinkurl{PaperIII.CanonicalTrianglePacking} makes \nolinkurl{PaperIII.nu3}
and \nolinkurl{PaperIII.nu3Star} the manuscript-facing quantities and
proves their equality with the edge-hypergraph implementations used by
AX1; it also proves \nolinkurl{PaperIII.tau3Star\ =\ PaperIII.nu3Star} and
identifies the cover-side AX1 interface with the packing-side statement
in Section 2. \nolinkurl{PaperIII.PaperImprovements} contains the
paper-facing certificates; \nolinkurl{PaperIII.PublicAPI},
\nolinkurl{PaperIII.OfPartition}, and \nolinkurl{PaperIII.Obstructions} expose
downstream interfaces. The directed canonical gate reports the following
axiom footprint for these theorem surfaces:

\begin{quote}
{\LeanSymbols\fontsize{8}{10}\selectfont
\noindent \nolinkurl{[propext,}\hspace{0.45em}\nolinkurl{Classical.choice,}\hspace{0.45em}\nolinkurl{Quot.sound]}\par
}
\end{quote}

The final assembly uses explicit bridges to the Haxell--Rödl/Yuster and
dense-decomposition formalizations. The cited asymptotic theorems remain
in the manuscript for mathematical provenance. The recorded author-side
public-root build completed 8,455 jobs with exit code 0; the subsequent
query-root build completed 8,444 jobs with exit code 0. That author-side
process was resumed after an application restart. Separately, the
external audit performed an uninterrupted reproduction from an absent
project build directory: the public root completed 8,455 jobs, the query
roots completed 8,444 jobs, and all eight axiom-query files exited
successfully. The immutable formal archive is sealed, and the internal
and external Lean gates pass within the scope stated in Section 13.

Appendix E gives a curated selection of paper-level and reusable
certificates. The complete implementation-level inventory remains in the
formal source rather than being duplicated in the manuscript.

\textbf{Table 2. Formalization surfaces.} Statuses in this table are
records of the sealed local freeze; they are not a substitute for
independent reproduction.

{\def\LTcaptype{none} % do not increment counter
\begin{longtable}[]{@{}
  >{\raggedright\arraybackslash}p{(\linewidth - 4\tabcolsep) * \real{0.3333}}
  >{\raggedright\arraybackslash}p{(\linewidth - 4\tabcolsep) * \real{0.3333}}
  >{\raggedright\arraybackslash}p{(\linewidth - 4\tabcolsep) * \real{0.3333}}@{}}
\toprule\noalign{}
\begin{minipage}[b]{\linewidth}\raggedright
Surface
\end{minipage} & \begin{minipage}[b]{\linewidth}\raggedright
Purpose
\end{minipage} & \begin{minipage}[b]{\linewidth}\raggedright
Current local status
\end{minipage} \\
\midrule\noalign{}
\endhead
\bottomrule\noalign{}
\endlastfoot
\nolinkurl{PaperIII} & aggregate/compatibility root & public-root build
completed 8,455 jobs; exit 0; imports the final theorem and public
API \\
\nolinkurl{PaperIII.MainNibble} & theorem and corollary assembly through
explicit AX1/AX2 interfaces & compiled in the recorded closure \\
\nolinkurl{PaperIII.Theorem_1_1_Final} & unconditional discharge of
AX1/AX2 and final theorem surfaces & compiled; final theorem footprint
is foundational-only \\
\nolinkurl{PaperIII.CanonicalTrianglePacking} & canonical \nolinkurl{nu3},
\nolinkurl{nu3Star}, \nolinkurl{tau3Star} and AX1 statement bridges &
compiled; exact target-side bridges recorded \\
\nolinkurl{PaperIII.CanonicalTrianglePackingGate} & statement and axiom
gate for the canonical interfaces & exit 0; foundational-only footprints
recorded \\
\nolinkurl{PaperIII.PaperImprovements} and
\nolinkurl{PaperIII.PaperImprovementsGate} & paper-facing corollaries and
selected auxiliary certificates & compiled in the consolidated gate
run \\
\nolinkurl{PaperIII.PublicAPI}, \nolinkurl{PaperIII.OfPartition},
\nolinkurl{PaperIII.Obstructions} & downstream API, split-partition
transport, and obstruction certificates & compiled in the recorded
closure \\
\nolinkurl{FREEZE_REPORT.md} and \nolinkurl{ESCAPE_HATCH_ASSESSMENT.md} &
immutable build, axiom and raw-token assessment & generated and included
in the sealed formal archive \\
\end{longtable}
}

At the \nolinkurl{SimpleGraph} level,
\nolinkurl{Theorem_1_1_of_splitPartition_uncond} accepts a finite
graph equipped with a clique/independent partition. The construction
\nolinkurl{SplitGraph.ofPartition} is connected to the original graph by
\nolinkurl{ofPartitionIso}; the vertex count, edge count, \(\nu_3\), and
\(\Phi\) are transported through graph-isomorphism invariance, including
\nolinkurl{nu3_congr_of_iso}. Other type-specialized theorem forms
remain in \nolinkurl{PaperIII.PublicAPI} for downstream use but are not
mathematical results of the paper and are not catalogued here.

The directed canonical-module compilation, public- and query-root
builds, eight compatible axiom-query files, escape-hatch assessment,
immutable source manifests, and internal audit pass their author-side
checks. The external audit independently reproduced the public and query
roots and verified the designated axiom surfaces against the same frozen
bytes.

No computation or regression script is a premise of the proof. The
exact-arithmetic and integer-programming checks remain supplementary
audits. The completed adversarial review and corpus-bounded prior-art
assessment are not substitutes for human peer review or a specialist
determination of priority; public tagging is a repository-level release
operation.

\subsection{11.7 Reusable formalization
components}\label{reusable-formalization-components}

The formalization also produced a standalone \nolinkurl{Contrib.*} library
whose modules import only Mathlib. Four groups are especially relevant
beyond this paper:

\begin{itemize}
\tightlist
\item
  finite packing/covering LP strong duality via a conic-Farkas argument
  (\nolinkurl{Contrib.LPStrongDuality});
\item
  variance-sensitive concentration and martingale tails
  (\nolinkurl{Contrib.BennettBernstein}, \nolinkurl{Contrib.HoeffdingUpper},
  \nolinkurl{Contrib.ConditionalHoeffding}, and
  \nolinkurl{Contrib.FreedmanBernstein});
\item
  Dirac-type Hamiltonicity and matching results
  (\nolinkurl{Contrib.DiracHamiltonian} and
  \nolinkurl{Contrib.SpreadMatchingDirac});
\item
  finite max-flow/min-cut and selected design-theoretic components
  (\nolinkurl{Contrib.MaxFlowMinCut} and \nolinkurl{Contrib.FanoSTS}).
\end{itemize}

These modules are reusable formalization byproducts, not premises of
Theorem 1.1. They are held for separate review and possible upstream
submission. The present manuscript makes no claim that their APIs or
individual statements are novel relative to the current Mathlib library;
that comparison requires a separate upstream-oriented literature and API
review. Appendix E.6 records only the components closest to the
mathematical proof.

\begin{center}\rule{0.5\linewidth}{0.5pt}\end{center}

\subsection{13. Reproducibility}\label{reproducibility}

The mathematical source is this manuscript. The manuscript-facing formal
surfaces are \nolinkurl{PaperIII.Theorem_1_1_Final},
\nolinkurl{PaperIII.CanonicalTrianglePacking}, \nolinkurl{PaperIII.PublicAPI},
\nolinkurl{PaperIII.OfPartition}, and \nolinkurl{PaperIII.Obstructions}.
\nolinkurl{PaperIII.PaperImprovementsGate} is the byproduct gate, and
\nolinkurl{PaperIII.CanonicalTrianglePackingGate} checks the exact
correspondence between the canonical packing parameters and the
Nibble/Yuster model used to discharge AX1.

The formal development is sealed as
\nolinkurl{PAPER_III_lean_v1.4_freeze.zip}, with SHA-256
\nolinkurl{79ee24c38fd776bc2585a0c3c996e30817f0829fc5064463bdbde0fa2d3d7104}.
It contains 704 Lean source files and the three pinned Lake/toolchain
files listed in the 707-entry source manifest. The recorded verification
uses Lean 4.28.0 and Mathlib v4.28.0 at commit
\nolinkurl{8f9d9cff6bd728b17a24e163c9402775d9e6a365}.

The available local verification record is summarized below.

\textbf{Table 3. Sealed local verification record.}

{\def\LTcaptype{none} % do not increment counter
\begin{longtable}[]{@{}
  >{\raggedright\arraybackslash}p{(\linewidth - 4\tabcolsep) * \real{0.3333}}
  >{\raggedright\arraybackslash}p{(\linewidth - 4\tabcolsep) * \real{0.3333}}
  >{\raggedright\arraybackslash}p{(\linewidth - 4\tabcolsep) * \real{0.3333}}@{}}
\toprule\noalign{}
\begin{minipage}[b]{\linewidth}\raggedright
Check
\end{minipage} & \begin{minipage}[b]{\linewidth}\raggedright
Log or command
\end{minipage} & \begin{minipage}[b]{\linewidth}\raggedright
Current result
\end{minipage} \\
\midrule\noalign{}
\endhead
\bottomrule\noalign{}
\endlastfoot
public aggregate-root build &
\nolinkurl{gate_logs/run_2026-08-22_CCS_NOTEBOOK456_resumed/logs/02_build_public_root_clean.log}
& exit 0; 8,455 jobs completed; clean-origin process resumed after an
application restart \\
query-root completion &
\nolinkurl{gate_logs/run_2026-08-22_CCS_NOTEBOOK456_resumed/logs/03_build_query_roots_incremental.log}
& exit 0; 8,444 jobs completed \\
canonical statement-and-axiom gate &
\nolinkurl{PaperIII.CanonicalTrianglePackingGate} & exit 0; target-side
integral and fractional bridges checked \\
complete theorem-level axiom record & eight \nolinkurl{FreezeAxioms*.lean}
files & exit 0; only \nolinkurl{propext}, \nolinkurl{Classical.choice}, and
\nolinkurl{Quot.sound} \\
quantitative Lean-closure regression &
\nolinkurl{FreezeAxiomsAuditClosure.lean} and
\nolinkurl{FreezeAxiomsAX1Closure.lean} & all named closure surfaces
present with foundational-only footprints \\
escape-hatch assessment & raw-token and focused build/footprint checks &
no active \nolinkurl{sorry} warning, no active \nolinkurl{native_decide}, and
no project axiom in a canonical surface \\
immutable seal & source/package manifests, ZIP and SHA-256 sidecar & 707
source entries; archive hash verified \\
\end{longtable}
}

The recorded verification used one unchanged source snapshot. Its build
sequence was

\begin{quote}
{\LeanSymbols\fontsize{8}{10}\selectfont
\noindent \nolinkurl{lake}\hspace{0.45em}\nolinkurl{build}\hspace{0.45em}\nolinkurl{PaperIII}\par
\noindent \nolinkurl{lake}\hspace{0.45em}\nolinkurl{build}\hspace{0.45em}\nolinkurl{BKLO.MainDenseUnconditional}\hspace{0.45em}\nolinkurl{Nibble.AX1Closed}\hspace{0.45em}\nolinkurl{PaperIII.CanonicalTrianglePacking}\hspace{0.45em}\nolinkurl{PaperIII.Obstructions}\hspace{0.45em}\nolinkurl{PaperIII.PaperImprovementsGate}\hspace{0.45em}\nolinkurl{PaperIII.PublicAPI}\hspace{0.45em}\nolinkurl{PaperIII.Theorem_1_1_Final}\par
\noindent \nolinkurl{#print}\hspace{0.45em}\nolinkurl{axioms}\hspace{0.45em}\nolinkurl{PaperIII.Theorem_1_1}\par
\noindent \nolinkurl{#print}\hspace{0.45em}\nolinkurl{axioms}\hspace{0.45em}\nolinkurl{PaperIII.AX1_holds}\par
\noindent \nolinkurl{#print}\hspace{0.45em}\nolinkurl{axioms}\hspace{0.45em}\nolinkurl{PaperIII.AX2_holds}\par
}
\end{quote}

The author-side project began with no \nolinkurl{.lake} directory or
compiled project object. The desktop application restarted during the
first public-root process; the unchanged project was then resumed
incrementally and completed successfully. The immutable freeze therefore
records \nolinkurl{PASS_CLEAN_ORIGIN_RESUMED}, not a single
uninterrupted author invocation. It preserves that disclosure, verbatim
output, component identifiers, source manifest, archive manifest, and
SHA-256 inventories. The external audit subsequently supplied an
uninterrupted reproduction from an absent project build directory, using
the pinned dependency revisions: 8,455 public-root jobs and 8,444
query-root jobs completed with exit code 0, and all eight axiom-query
files passed. Those declaration-level checks establish build conformance
and axiom footprints; they do not replace independent mathematical
rederivation.

Two archived legacy modules, \nolinkurl{Ax2/PartB/Axioms.lean} and
\nolinkurl{Ax2/PartA/Wlog.lean}, contain project-local axiom declarations.
No source file imports either module, and neither declaration occurs in
the dependency footprint of a canonical release surface. Accordingly,
the precise formal claim is declaration-level: the canonical theorem and
release interfaces have no project-local mathematical axiom; the
complete historical source tree is not claimed to contain no
\nolinkurl{axiom} token.

The proof itself remains in the manuscript. Regression scripts and
computational audits are not hypotheses, reductions, or proof steps.

The canonical repository package contains the complete Lean source,
toolchain, manifest, build and axiom logs, manuscript sources, and
SHA-256 inventories. The public tag and commit bind those files at
repository level; no commit identifier is asserted inside the manuscript
itself.

Written materials are intended for distribution under CC BY-NC 4.0.
Author-owned verification and regression code is intended for
distribution under the PolyForm Noncommercial License 1.0.0; third-party
dependencies retain their own licenses.

\textbf{Review boundary:} the independent reproduction and external
adversarial audit are complete, and the novelty review is
corpus-bounded. Human peer review and specialist priority determination
have not been performed. The public tag and commit are repository
metadata rather than premises of the theorem.

\begin{center}\rule{0.5\linewidth}{0.5pt}\end{center}

\subsection{Appendix C. Computational
Audits}\label{appendix-c.-computational-audits}

\subsection{C.1 Exact common-profile
LP}\label{c.1-exact-common-profile-lp}

\nolinkurl{verify_common_profile_lp.py} enumerates the vertices of the
symmetrized cover polyhedron using exact rational Gaussian elimination
and compares the optimum with (3.5).

\subsection{C.2 Exact fractional
margin}\label{c.2-exact-fractional-margin}

\nolinkurl{verify_fractional_margin.py} checks Theorem 4.2 for

\[
3\le p\le80,
\qquad
1\le q\le2p,
\qquad
0\le d\le p
\]

using exact rational arithmetic.

\subsection{C.3 Small ILP audits}\label{c.3-small-ilp-audits}

\nolinkurl{verify_factor_rounding.py} and
\nolinkurl{verify_shifted_center.py} compute exact integral triangle
packings on small instances using a binary edge-capacity ILP. They
verify the stated upper bounds on \(\Phi(G)\).

\subsection{C.4 Polarization}\label{c.4-polarization}

\nolinkurl{verify_polarization.py} performs exhaustive small-profile and
randomized exact-integer checks of Lemma 6.1.

\subsection{C.5 Divisibility}\label{c.5-divisibility}

\nolinkurl{verify_divisibility_correction.py} exhaustively verifies the
path parity construction through order eighteen, checks the fixed
minimum-degree threshold exactly, and follows the full parity and
modulo-three correction on randomized dense residuals.

\begin{center}\rule{0.5\linewidth}{0.5pt}\end{center}

\subsubsection{\texorpdfstring{Remark 2.3\(^{\prime}\) --- Necessity of
a density
hypothesis}{Remark 2.3\^{}\{\textbackslash prime\} --- Necessity of a density hypothesis}}\label{remark-2.3prime-necessity-of-a-density-hypothesis}

\textbf{Lean certificates.}
\nolinkurl{PaperIII.ax2_divisibility_degree_insufficient},
\nolinkurl{PaperIII.ax2_density_necessary_K7_minus_two_triangles},
re-exported through \nolinkurl{PaperIII.Obstructions}.

\subsubsection{Corollary 10.2b --- Sharp quadratic term and forced
linear
coefficient}\label{corollary-10.2b-sharp-quadratic-term-and-forced-linear-coefficient}

\textbf{Lean certificate.} \nolinkurl{PaperIII.Corollary_1_2_sharp},
assembled from \nolinkurl{PaperIII.Corollary_1_2},
\nolinkurl{PaperIII.Byproduct_leading_constant_forced}, and
\nolinkurl{PaperIII.Byproduct_completeSplit_cp_sharp}. The declaration
\nolinkurl{Byproduct_leading_constant_forced} certifies the assertion
about the linear coefficient; sharpness of the quadratic coefficient
follows from the exact complete-split value as \(p\to\infty\).

\subsubsection{Corollary E.1.1 --- Single factor-assignment
packing}\label{corollary-e.1.1-single-factor-assignment-packing}

\textbf{Lean certificate.}
\nolinkurl{Byproduct_factorization_assignment_packing}.

\subsubsection{Corollary E.1.2 --- Double-factorization lower
bound}\label{corollary-e.1.2-double-factorization-lower-bound}

\textbf{Lean certificate.}
\nolinkurl{Byproduct_double_factorization_packing}.

\subsubsection{Corollary E.1.3 --- Reserved-gain packing
bound}\label{corollary-e.1.3-reserved-gain-packing-bound}

\textbf{Lean certificate.}
\nolinkurl{Byproduct_reserved_gain_packing_bound_subset}.

\subsubsection{\texorpdfstring{Corollary E.2.1 --- Short-corridor
constant
\(3/2\)}{Corollary E.2.1 --- Short-corridor constant 3/2}}\label{corollary-e.2.1-short-corridor-constant-32}

\textbf{Lean certificate.}
\nolinkurl{Byproduct_short_corridor_constant_three_halves}.

\subsubsection{Corollary E.2.2 --- Packing-to-partition
bridge}\label{corollary-e.2.2-packing-to-partition-bridge}

\textbf{Lean certificates.} \nolinkurl{Byproduct_cp_le_Phi},
\nolinkurl{Byproduct_low_corridor_cp_three_halves},
\nolinkurl{Byproduct_effective_low_corridor_cp_bound}.

\subsubsection{Corollary E.3.1 --- Exact fractional
optimum}\label{corollary-e.3.1-exact-fractional-optimum}

\textbf{Lean certificate.}
\nolinkurl{Byproduct_commonProfile_tau3Star_eq_F}.

\subsubsection{Corollary E.3.2 --- Phase
regions}\label{corollary-e.3.2-phase-regions}

\textbf{Lean certificates.}
\nolinkurl{Byproduct_commonProfile_tau3Star_eq_uniform_region},
\nolinkurl{Byproduct_commonProfile_tau3Star_eq_separated_region},
\nolinkurl{Byproduct_commonProfile_tau3Star_eq_hot_region}.

\subsubsection{Corollary E.4.1 --- Fractional slack lower
bound}\label{corollary-e.4.1-fractional-slack-lower-bound}

\textbf{Lean certificate.}
\nolinkurl{Byproduct_unified_fractional_margin}.

\subsubsection{Corollary E.5.1 --- Exact value and normal
form}\label{corollary-e.5.1-exact-value-and-normal-form}

\textbf{Lean certificates.}
\nolinkurl{Byproduct_completeSplit_cp_exact_value},
\nolinkurl{Byproduct_completeSplit_cp_sharp}.

\subsubsection{Corollary E.5.2 --- Forced linear
coefficient}\label{corollary-e.5.2-forced-linear-coefficient}

\textbf{Lean certificates.}
\nolinkurl{Byproduct_leading_constant_forced}, packaged with the upper
bound and exact family as \nolinkurl{Corollary_1_2_sharp}.

\subsubsection{E.6.1 Finite packing/covering LP
duality}\label{e.6.1-finite-packingcovering-lp-duality}

\textbf{Lean certificates.} \nolinkurl{Byproduct_triangle_LP_duality};
standalone form \nolinkurl{Contrib.LPDuality.lp_strong_duality} in
\nolinkurl{Contrib.LPStrongDuality}.

\subsubsection{E.6.2 Dirac-type matching and dense
correction}\label{e.6.2-dirac-type-matching-and-dense-correction}

\textbf{Lean certificates.}
\nolinkurl{Byproduct_dirac_perfect_matching_even},
\nolinkurl{Byproduct_dirac_near_perfect_matching},
\nolinkurl{Byproduct_dense_divisible_correction_edges}; standalone
matching form in \nolinkurl{Contrib.SpreadMatchingDirac}.

\subsubsection{E.6.3 Uniform nibble
interface}\label{e.6.3-uniform-nibble-interface}

\textbf{Lean certificate.}
\nolinkurl{Byproduct_almostPerfectMatching_uniform}.

\subsubsection{Corollary E.7.1 --- Attainment and fractional edge
bound}\label{corollary-e.7.1-attainment-and-fractional-edge-bound}

\textbf{Lean certificates.} \nolinkurl{Byproduct_nu3_achieved},
\nolinkurl{Byproduct_nu3star_le_edges_div_three}.

\subsubsection{Corollary E.7.2 --- Majority-nibble edge-gap
form}\label{corollary-e.7.2-majority-nibble-edge-gap-form}

\textbf{Lean certificate.}
\nolinkurl{Byproduct_nu3star_sub_nu3_le_edge_gap_most}.

\subsubsection{Appendix E. Selected Formalization-Derived
Corollaries}\label{appendix-e.-selected-formalization-derived-corollaries}

The theorem names follow \nolinkurl{PaperIII.PaperImprovements} and the
standalone \nolinkurl{Contrib.*} package. Their release status is governed
by the consolidated build and axiom gate in Section 13.

\subsubsection{Corollary E.5.2 --- Forced linear
coefficient}\label{corollary-e.5.2-forced-linear-coefficient-1}

The weaker certificate \nolinkurl{Byproduct_leading_constant_sharp},
which follows immediately from the exact normal form, is not separately
catalogued.

\subsubsection{E.6.3 Uniform nibble
interface}\label{e.6.3-uniform-nibble-interface-1}

The accompanying standalone library also contains concentration,
Hamiltonicity, max-flow/min-cut, codegree-counting, hypergeometric-tail,
and design-theoretic modules. Section 11.7 records their scope without
presenting them as additional corollaries of this paper.

\subsubsection{Corollary 10.2a --- Exact complete-split
benchmark}\label{corollary-10.2a-exact-complete-split-benchmark}

The Lean development records the upper bound, lower bound, and equality
as \nolinkurl{Byproduct_completeSplit_cp_le_exact_value},
\nolinkurl{Byproduct_completeSplit_cp_ge_exact_value}, and
\nolinkurl{Byproduct_completeSplit_cp_exact_value}. Appendix E.5
retains only the exact normal form and the logically distinct
forced-linear-coefficient statement.

\subsubsection{Packing-form corollaries}\label{packing-form-corollaries}

The following corollaries restate results from Sections 5 and 7 directly
in packing language. They are formalized as
\nolinkurl{factorization_assignment_packing},
\nolinkurl{double_factorization_packing}, and
\nolinkurl{reserved_gain_packing_bound_subset}. All three are
unconditional: none uses Theorem 2.1 or Theorem 2.3.

\subsection{References}\label{references}

{[}1{]} P. Allen, J. Böttcher, O. Cooley, and R. Mycroft, ``Tight cycles
and regular slices in dense hypergraphs,'' \emph{Journal of
Combinatorial Theory, Series A} \textbf{149} (2017), 30--100.

{[}2{]} B. Barber, D. Kühn, A. Lo, and D. Osthus, ``Edge-decompositions
of graphs with high minimum degree,'' \emph{Advances in Mathematics}
\textbf{288} (2016), 337--385.

{[}3{]} T. F. Bloom, ``Erdős Problem \#81,'' \emph{Erdős Problems},
\url{https://www.erdosproblems.com/81}, accessed July 6, 2026.

{[}4{]} O. V. Borodin, A. V. Kostochka, and D. R. Woodall, ``List edge
and list total colourings of multigraphs,'' \emph{Journal of
Combinatorial Theory, Series B} \textbf{71} (1997), no. 2, 184--204.

{[}5{]} G.-T. Chen, P. Erdős, and E. T. Ordman, ``Clique partitions of
split graphs,'' in \emph{Combinatorics, Graph Theory, Algorithms and
Applications} (Beijing, 1993), World Scientific, 1994, pp.~21--30.

{[}6{]} G. A. Dirac, ``Some theorems on abstract graphs,''
\emph{Proceedings of the London Mathematical Society} (3) \textbf{2}
(1952), 69--81.

{[}7{]} F. Dross, ``Fractional triangle decompositions in graphs with
large minimum degree,'' \emph{SIAM Journal on Discrete Mathematics}
\textbf{30} (2016), no. 1, 36--42.

{[}8{]} P. Erdős, E. T. Ordman, and Y. Zalcstein, ``Clique partitions of
chordal graphs,'' \emph{Combinatorics, Probability and Computing}
\textbf{2} (1993), no. 4, 409--415.

{[}9{]} D. Gale and L. S. Shapley, ``College admissions and the
stability of marriage,'' \emph{The American Mathematical Monthly}
\textbf{69} (1962), no. 1, 9--15.

{[}10{]} F. Galvin, ``The list chromatic index of a bipartite
multigraph,'' \emph{Journal of Combinatorial Theory, Series B}
\textbf{63} (1995), no. 1, 153--158.

{[}11{]} P. E. Haxell and V. Rödl, ``Integer and fractional packings in
dense graphs,'' \emph{Combinatorica} \textbf{21} (2001), 13--38.

{[}12{]} P. Keevash, ``The existence of designs,'' arXiv:1401.3665,
revised 2018.

{[}13{]} D. König, ``Über Graphen und ihre Anwendung auf
Determinantentheorie und Mengenlehre,'' \emph{Mathematische Annalen}
\textbf{77} (1916), 453--465.

{[}14{]} M. Krivelevich, ``On a conjecture of Tuza about packing and
covering of triangles,'' \emph{Discrete Mathematics} \textbf{142}
(1995), 281--286.

{[}15{]} J. P. Traverso Gianini, ``Affine profile reduction for
fractional triangle packings in split graphs'' (Paper I in the series),
preprint, July 2026.

{[}16{]} J. P. Traverso Gianini, ``Complete-split extremizers for a
fractional triangle-cover functional on chordal graphs'' (Paper II in
the series), preprint, July 2026.

{[}17{]} R. Yuster, ``Integer and fractional packing of families of
graphs,'' \emph{Random Structures \& Algorithms} \textbf{26} (2005),
110--118.

\end{document}
